\subsection{Model Transfer from Shake Flask to Bioreactor and Temperature Downshift}
\subsubsection{State Estimation Across Scale and Temperature Shift}
\rev{The shake flask culture at 36.5\,°C (Dataset 1, grey) provided the base mechanistic model, which was then adapted to bioreactor cultures at the same temperature (Dataset 2, purple) and following a temperature downshift to 32\,°C (Dataset 3, teal).} The resulting EnKF predictions for all three conditions are shown in Figure \ref{fig:T127_overlay}. The filter facilitates recursive parameter and state updates throughout the culture without model reparametrisation.

The higher cell growth and mAb productivity observed at bioreactor scale were both well captured. The EnKF accurately reproduced the overall glucose consumption trends across all three conditions. A slight misalignment was observed under the 32°C condition, which may be attributed to uncertainty in the composition of the proprietary Feed X added in addition to the commercial Feed C. Although Feed X was modelled as part of the nutrient input, its compositional uncertainty was not explicitly represented in the EnKF, meaning its variability was not propagated through the ensemble. Ammonia accumulation plateaued earlier at bioreactor scale compared to shake flask, with the EnKF recovering the general trajectory. Importantly, the filter did not track the ammonia dip between days 6 and 10 for Bioreactor 36.5°C dataset, indicating that it did not overfit to what may represent measurement noise rather than a true metabolic shift.


Following the temperature downshift, glutamine concentrations in the Bioreactor 32°C dataset were substantially elevated, and this was reliably captures by the filter. A distinct shift in lactate metabolism was also observed in response to the temperature downshift, with net lactate consumption occurring after day 6, in contrast to the accumulation observed in shake flask and 36.5°C bioreactor conditions. The lactate consumption trend is successfully captured but the end concentration is slightly lower than that described by the EnKF. Glutamate concentrations varied across datasets, yet the predicted trajectories followed the observed profiles across all three conditions. Asparagine dynamics exhibited clear scale dependence, with the EnKF accurately describing the rate and extent of depletion.

\begin{figure}[H]
    \centering
    \includegraphics[width=0.95\textwidth]{Figs/T127_overlay_errorbars.png}
    \caption{\rev{EnKF-based state estimation for Cell Line A across scale and temperature downshift. Trajectories are shown for shake flask at 36.5\,°C (Dataset 1, grey; base model calibration), bioreactor at 36.5\,°C (Dataset 2, purple), and bioreactor with temperature downshift to 32\,°C (Dataset 3, teal). Solid lines show the EnKF posterior mean; markers show experimental measurements ($\pm$standard deviation).}}
    \label{fig:T127_overlay}
\end{figure}

\subsubsection{Parameter Evolution for Knowledge Transfer Across Scale and Temperature Shift}

\rev{Figure~\ref{fig:T127_para} shows the time evolution of selected model parameters for Cell Line A across scale and temperature conditions, illustrating how the EnKF progressively adapts the base model to each new operating condition.}

\paragraph{Cell growth and productivity}

Parameter values remain relatively constant across culture duration in the shake flask experiments. This is expected given that this dataset was used for model parametrisation and the initial ensemble adequately captures baseline behaviour. The trajectory of the maximum specific growth rate $\mu_{max}$ and death rate $\mu_{d,max}$ aligns well with the observed viable cell density profile shown in Figure~\ref{fig:T127_overlay} (a). In the bioreactor experiments conducted at 36.5°C, $\mu_{max}$ is increased during early exponential growth, reflecting the higher viable cell densities achieved in this controlled environment with stable pH and dissolved oxygen tension. Interestingly,  the trajectory of $\mu_{d,max}$ shows higher rate of cell death in bioreactors compared to shake flask cultures, potentially because the higher viable cell density achieved. In the 32°C bioreactor, however, a clear reduction in $\mu_{d,max}$ is observed following the temperature downshift, indicating delayed and slower cell death. This behaviour is consistent with the expected benefits of mild hypothermia in CHO cultures, which has been shown to prolong viability and extend the productive phase \cite{Sou2015HowGlycosylation, Sou2017Model-basedHypothermia}. These advantages are also evident in the trajectory of the titre production parameter $m_{mAb}$, which reaches its highest values under 32°C, reflecting enhanced specific productivity with temperature downshift. 

\begin{figure}[H]
    \centering
    \includegraphics[width=0.95\textwidth]{Figs/para_compare_scale.png}
    \caption{\rev{Time evolution of eight estimated model parameters across scale and temperature downshift for Cell Line A, shown for shake flask (Dataset 1, grey), bioreactor at 36.5\,°C (Dataset 2, purple), and bioreactor with temperature downshift to 32\,°C (Dataset 3, teal). Parameters shown are the maximum specific growth rate ($\mu_{\max}$), maximum specific death rate ($\mu_{\text{death,max}}$), non-growth associated mAb production rate ($m_{\text{mAb}}$), yield of biomass on glutamine ($Y_{X/Gln}$), yield of biomass on ammonia ($Y_{X/Amm}$), yield of biomass on lactate ($Y_{X/Lac}$), and lactate consumption activation constants ($Lac_{\max1}$, $Lac_{\max2}$). Solid lines show the posterior mean; shaded bands show the ensemble spread.}}
    \label{fig:T127_para}
\end{figure}

\paragraph{Scale-dependent glutamine utilization}

The yield of biomass from glutamine $Y_{x,Gln}$ remains relatively constant with time in the shake flasks but increases in both bioreactor cultures from day 6 onwards, reflecting enhanced glutamine utilisation for cell growth at larger scale and in a controlled environment. In GS-CHO cell lines, glutamine is synthesised from glutamate and ammonia via an ATP-dependent glutamine synthetase reaction \cite{Lin2019AttenuatedBiologics}. As ATP is primarily generated through oxidative phosphorylation, the rate of this conversion is sensitive to oxygen availability \cite{Wilson2017OxidativeMetabolism}. Improved oxygen transfer in bioreactors supports sustained ATP production, thereby enabling more efficient glutamine synthesis from ammonia. This mechanistic interpretation is further supported by a parallel increase in the yield of biomass from ammonia $Y_{x,Amm}$. Although glutamine concentration profiles appear similar in the shake flask and 36.5°C bioreactor until day 12 as seen in Figure \ref{fig:T127_overlay} (e), the dynamic trajectory of $Y_{x,Gln}$ inferred by the EnKF reveals scale-dependent effects as early as day 6. 


\paragraph{Lactate consumption induced by temperature shifts}

The yield of biomass on lactate $Y_{X,Lac}$ reaches its highest values in the 32°C bioreactor, suggesting increased utilisation of lactate as an energy source following the temperature downshift. Scaling up from shake flask to bioreactor at 36.5°C does not substantially alter lactate accumulation, and this is reflected in the relatively stable trajectories of $Lac_{max1}$ and $Lac_{max2}$, which remain close to the base model. These parameters represent kinetic constants that regulate the activation of lactate consumption pathways. In the 32°C cultures, $Lac_{max1}$ increases during the later stages indicating lactate accumulation, while $Lac_{max2}$ decreases sharply after the downshift due to net lactate consumption. The parameter trajectories provide evidence for the stronger influence of $Lac_{max2}$ on net lactate consumption, reflecting the bidirectional nature of lactate metabolism.


\subsection{Model Transfer from Cell-line A to B under Different Feeds}

\subsubsection{State Estimation Across Cell-lines and Feeding Conditions}

The base model, originally parametrised using Cell Line A shake flask data supplemented with Feed C, was transferred to the Cell Line B. Three feeding strategies were evaluated, Feed C (Dataset 4, orange), Feed U (Dataset 5, green), and Feed U plus 40\% (Dataset 6, blue), with all studies conducted in shake flasks. Despite differences in cell line and feed compositions, the EnKF successfully adapted the base model to capture cell growth, productivity, and metabolite dynamics for all three cultures, as shown in Figure \ref{fig:GS46_overlay}.


\begin{figure}[H]
    \centering
    \includegraphics[width=0.95\textwidth]{Figs/gs46_overlay_errorbars.png}
    \caption{\rev{EnKF-based state estimation for Cell Line B across feeding strategies. Trajectories are shown for Feed C (Dataset 4, orange), Feed U (Dataset 5, green), and Feed U plus 40\% (Dataset 6, navy), all conducted as shake flask cultures at 36.5\,°C. Solid lines show the EnKF posterior mean; markers show experimental measurements ($\pm$standard deviation). All predictions are initialised from the Cell Line A base model (Dataset 1).}}
    \label{fig:GS46_overlay}
\end{figure}



Distinct growth and death profiles were observed across the feeding strategies. Compared to Feed C, Feed U and Feed U plus 40\% both resulted in a steeper decline in viable cell density following peak growth. These differences were accurately captured by the EnKF, even though all predictions were based on the same underlying model structure. The mAb titre profiles were also well reproduced, with Feed C yielding the highest final titre of approximately 2.5 g L\textsuperscript{-1}. This represents a nearly four-fold increase compared to the baseline model calibrated on Cell Line A, yet the EnKF successfully captured this shift, demonstrating transferability across systems with markedly different productivity levels.

The higher concentrations of glucose and asparagine in Feed U and Feed U plus 40\% resulted in different metabolite profiles compared to Feed C. In particular, excess asparagine feeding led to substantial ammonia accumulation, with concentrations exceeding 20mM by the end of the culture under Feed U plus 40\%. While the EnKF described the general trend well, some underprediction was observed during the final stage in this condition. During filter tuning, it was found that a higher ammonia measurement noise was required to stabilise the EnKF in the Feed U plus 40\% case. Reducing the measurement noise resulted in a closer fit to the ammonia data but introduced instability at later time points, as the filter was forced too closely towards the measurements and propagated aggressive corrections into correlated states, particularly glutamine, glutamate and asparagine. The slight mismatch observed is therefore expected, given the greater uncertainty in ammonia measurements at high concentrations. Nonetheless, the EnKF well represented the broader dynamics of nitrogen metabolism across conditions, including elevated glutamine and glutamate concentrations in both Feed U and Feed U plus 40\%.

Lactate profiles further highlight the flexibility of the EnKF framework. The base model parametrised on the Cell Line A shake flask dataset exhibited net lactate accumulation under Feed C, as shown in Figure \ref{fig:T127_overlay} (f). In contrast, Cell Line B culture under the same scale and feeding strategy showed net lactate consumption, indicating this behaviour is cell line-specific. Feed U and Feed U plus 40\% resulted in even stronger lactate consumption, with near-complete depletion by the end of the culture in the latter case. These distinct profiles were accurately described using the same underlying mechanistic model, made possible by the dynamic evolution of the parameter ensemble. Such flexibility would not be attainable under a fixed-parameter model calibration. {\color{blue} Crucially, this adaptability is not achieved at the cost of predictive accuracy. EnKF predictions for Cell Line B under Feed C were benchmarked against a reparametrised open-loop simulation following Kotidis \textit{et al.}~\cite{Kotidis2019UseCultures}, which first performs Global Sensitivity Analysis (GSA) to identify the most influential model parameters and then re-fits them specifically for target dataset. This constitutes a rigorous gold standard, but one that is both labour-intensive and requires considerable modelling expertise to execute for every new system. The EnKF requires neither: no prior sensitivity analysis campaign is needed, and parameter sensitivity emerges implicitly through ensemble spread reduction as measurements are assimilated. Despite beginning from Cell Line A base parameters without dataset-specific fitting, the EnKF achieved comparable RMSE and R\textsuperscript{2} across all state variables, demonstrating that sequential measurement assimilation can substitute for a priori reparametrisation. Full quantitative metrics are provided in Supplementary Information~\ref{SI:comparison}.}

\subsubsection{Parameter Evolution for Knowledge Transfer Across Cell-lines and Feeding Conditions}

\rev{Figure~\ref{fig:GS46_para} shows the time evolution of selected model parameters for Cell Line B across feeding strategies, revealing how the EnKF differentiates feed-specific metabolic behaviour.}

\paragraph{Cell growth and productivity} The trajectory of the maximum specific death rate $\mu_{d,max}$ provides a clear distinction between the commercial Feed C and the in-house formulations, Feed U and Feed U plus 40\%. Ensemble-based uncertainty narrows after the exponential growth phase, reflecting the higher sensitivity of $\mu_{d,max}$ during later stages of cell culture. In Cell Line B supplemented with Feed C, $\mu_{d,\text{max}}$ remains relatively constant throughout the culture, within the uncertainty of the parameter estimate. In contrast, Feed U shows an increase in the mean death rate during the stationary phase, which is expected given its limited nutrient composition compared to the commercial feed and the associated rise in ammonia levels. Feed U plus 40\% exhibits an even higher $\mu_{d,\text{max}}$, consistent with the pronounced ammonia accumulation observed in this condition and its strong correlation with increased cell death, as illustrated in Figure~\ref{fig:GS46_overlay} (a). Among the three feeding strategies, Feed C supports the highest titre yields, while the two in-house formulations result in comparable but lower final titres, as shown in Figure \ref{fig:GS46_overlay} (b). This suggests that improved cell viability under Feed C is accompanied by higher non-growth-associated productivity $m_{mAb}$, contributing to enhanced overall production.


\begin{figure}[H]
    \centering
    \includegraphics[width=0.95\textwidth]{Figs/para_compare_feed.png}
    \caption{\rev{Time evolution of eight estimated model parameters across feeding strategies for Cell Line B, shown for Feed C (Dataset 4, orange), Feed U (Dataset 5, green), and Feed U plus 40\% (Dataset 6, navy). Parameters shown are the maximum specific death rate ($\mu_{\text{death,max}}$), non-growth associated mAb production rate ($m_{\text{mAb}}$), yield of biomass on glucose ($Y_{X/Glc}$), yield of biomass on asparagine ($Y_{X/Asn}$), ammonia inhibition constant ($KI_{\text{Amm}}$), yield of biomass on lactate ($Y_{X/Lac}$), and lactate consumption activation constants ($Lac_{\max1}$, $Lac_{\max2}$). Solid lines show the posterior mean; shaded bands show the ensemble spread.}}
    \label{fig:GS46_para}
\end{figure}

\paragraph{Effect of excess feeding on biomass yield and nitrogen metabolism}

To assess how elevated glucose and asparagine in the in-house feeds influence growth in the new CHO-GS cell line, the temporal evolution of biomass yields from glucose $Y_{X,Glc}$ and asparagine $Y_{X,Asn}$ was examined. As shown in Figure~\ref{fig:GS46_para} (c), $Y_{X,Glc}$ exhibits broad ensemble spread for all conditions, with only modest changes in the parameter mean over time. This indicates that $Y_{X,Glc}$ is relatively insensitive to variations in glucose availability and does not strongly influence growth predictions.

In contrast, the $Y_{X,Asn}$ ensemble spread reduces rapidly, indicating that the filter gains confidence quickly about the contribution of asparagine to biomass formation. $Y_{X,Asn}$ is initially highest in Feed U plus 40\%, reflecting its elevated concentration. Over time, this parameter declines in Feed U with reduced uncertainty, indicating growing model confidence that asparagine contributes less to biomass. Interestingly, $Y_{X,Asn}$ in Feed U plus 40\% drops below that of Feed C by the end, despite significantly higher asparagine availability and uptake rate. This disparity suggests rerouting of excess asparagine through alternative nitrogen pathways. Previous studies have proposed that, under nitrogen-rich conditions, asparagine synthetase may catalyse a reaction analogous to asparaginase, converting asparagine and glutamate into aspartate and glutamine \cite{Quek2010MetabolicCulture}. While not conventionally reported in mammalian systems, this reaction has been included in CHO metabolic reconstructions to account for elevated ammonia accumulation under nitrogen-rich conditions \cite{Carinhas2013MetabolicTransition, Kyriakopoulos2014ACulture}. The parameter trajectory of $Y_{X,Asn}$ thus offers indirect evidence of such metabolic shifts, revealing insight not captured by concentration data alone.

To further investigate the impact of excess asparagine feeding, the ammonia inhibition parameter $K_{I,Amm}$ was evaluated. In Feed C, this parameter remains near its initial value with broad uncertainty, consistent with its use in model calibration. In contrast, both in-house feeds exhibit a sharp increase in $K_{I,Amm}$ entering stationary phase, accompanied by a rapid drop in uncertainty. This trend aligns with the timing of ammonia accumulation and declining $Y_{X,Asn}$, indicating that excess nitrogen progressively limits growth in later stages.      
\paragraph{Understanding lactate metabolism beyond concentration profiles}

To understand how lactate metabolism differs across feeding strategies, the biomass yield on lactate $Y_{X,Lac}$ and two kinetic parameters $Lac_{max1}$ and $Lac_{max2}$ were analysed. While concentration profiles confirm net lactate consumption under all conditions, they cannot distinguish whether this uptake supports biomass growth or other cellular functions.The parameter trajectories provide a clearer picture. All three conditions show elevated $Y_{X,Lac}$ early in culture, suggesting enhanced growth-coupled lactate utilisation in Cell Line B. However, only Feed C maintains high yields in later stages, while both in-house formulations show a decline. This supports the hypothesis that lactate may be increasingly rerouted toward non-growth-associated functions such as redox balancing under nutrient-rich conditions \cite{Morrissey2025NicotinamideCulture}.

{\color{blue}The trajectories of the lactate activation constants $Lac_{\max1}$ and $Lac_{\max2}$ further elucidate the dynamics of consumption across feeding strategies.} $Lac_{max2}$ declines across all feeds but most sharply under Feed U plus 40\%, in line with complete lactate depletion. In the same condition, $Lac_{max1}$ also drops steadily, indicating sustained consumption. In contrast, under Feed C and Feed U, $Lac_{max1}$ rises again, suggesting some lactate accumulation remains despite overall uptake. {\color{blue}Importantly,} although the concentration profiles indicate net lactate accumulation throughout the early culture phase, the parameter trajectories {\color{blue}reveal an earlier} underlying metabolic shift. Both $Lac_{\text{max1}}$ and $Lac_{\text{max2}}$ begin to decline as early as day 4, well before glucose or amino acid depletion. This behaviour mirrors previous observations in Cell Line B batch cultures \cite{Kyriakopoulos2013ComparativeProductivity} and supports literature findings that lactate consumption can also be triggered by high extracellular lactate rather than nutrient limitation \cite{Zagari2013HighCulture, Mulukutla2012OnCells}. {\color{blue}The converged $Lac_{\max1}$ and $Lac_{\max2}$ values for Cell Line B Feed C are consistent with those reported by the GSA guided reparametrised model for this specific dataset \cite{Kotidis2019UseCultures}, providing cross-validation of the EnKF parameter estimates.} \rev{The evolution of $Lac_{\max1}$ and $Lac_{\max2}$ across datasets also exemplifies a broader consideration in ensemble-based parameter estimation. Parameters that are jointly informed by available measurements may exhibit posterior correlation, such that their individual contributions cannot be fully resolved independently from data alone. Examining posterior correlations alongside ensemble spread therefore provides a more complete picture of parameter identifiability. A full correlation analysis across all six datasets is provided in Supplementary Information~\ref{SI:param_corr}.}


\subsection{\textcolor{blue}{EnKF State Estimation Under Data-Limited Scenarios}}

\subsubsection{\textcolor{blue}{Long-Term Prediction with Real-Time Measurement Updates}}
    \textcolor{blue}{As a culture progresses, measurements become available sequentially and only data collected up to the current time point can be used. At each measurement occasion, the EnKF updates the parameter ensemble and launches an open-loop simulation forward to harvest — a forward prediction that uses the inferred parameters but incorporates no further measurement updates. As more data are assimilated, parameter estimates improve and forecasts progressively converge toward the true culture trajectory.} Figure \ref{fig:long-term} illustrates the long-term prediction performance of EnKF when applied to Cell Line B cultured with Feed U plus 40\%, a system that differs significantly from the Cell Line A baseline both in cell line and feeding strategy. By day 7, EnKF predictions already substantially improved upon the baseline mechanistic model for most states, including glucose, glutamine, lactate, and glutamate. By day 9, accurate forecasts are achieved for glucose, glutamine, glutamate, and asparagine. By day 12, all state trajectories align closely with those obtained using the full dataset, confirming the framework’s capacity to refine harvest-day outcomes in real time. The corresponding long-term prediction results for the remaining datasets are provided in Supplementary Information \ref{long-term}. Such progressively improving forecasts can support model-predictive control, where the updated parameter ensemble informs the optimisation of future control actions using the most recent process knowledge.
    However, the timing at which accurate long-term prediction becomes possible is case-dependent. It reflects both the degree of deviation from the base model and the availability of information on key metabolic shifts. For instance, the sharp decline in viable cell density near the end of culture is not reflected in predictions from day 7 or day 9, as this behaviour had not yet emerged in the available data. Since the EnKF builds upon a mechanistic model originally developed for a different system, its ability to adapt relies on the progressive assimilation of system-specific measurements. Thus, while the framework is powerful in adapting to new conditions, its predictive accuracy is fundamentally determined by the timing and availability of incoming data.

\begin{figure}[H]
    \centering
    \includegraphics[width=0.95\textwidth]{Figs/long_term_pred_CHO_GS46_F_all_pl40.png}
    \caption{\textcolor{blue}{Long-term EnKF predictions of Cell Line B Feed U plus 40\% initiated from day 7 (dashed), day 9 (dotted), and day 12 (dash-dot), compared against experimental measurements and the base Kotidis \textit{et al.}~\cite{Kotidis2019Model-basedCulture} model (solid red). Forecasts are open-loop simulations using the parameter ensemble updated up to each initiation day, with predictions progressively improving over the base model as measurements are assimilated.}}
    \label{fig:long-term}
\end{figure}
\subsubsection{\textcolor{blue}{State Estimation Under Irregular Measurement Schedules}}
{\color{blue}
    In practical bioprocess monitoring, measurements are not always collected at uniform intervals. Equipment downtime, operational constraints, or deliberate sampling strategies can result in irregular or missing time points. To evaluate the robustness of the EnKF framework under such conditions, a scenario was constructed in which measurements are available only at alternating 48-hour and 72-hour intervals, simulating a realistically sparse and non-uniform sampling schedule. Figure \ref{fig:irregular} shows the EnKF state estimation results for Cell Line B Feed U plus 40\% under this irregular schedule. Despite the reduced measurement frequency and non-uniform spacing, the EnKF successfully tracks all eight state variables throughout the culture, maintaining close agreement with the available experimental data. This demonstrates that the Bayesian update mechanism of the EnKF is inherently robust to irregular measurement schedules: each update simply incorporates the next available observation regardless of when it arrives, without requiring any structural changes to the algorithm. The corresponding results for all remaining datasets are provided in Supplementary Information \ref{SI:irregular}.}

\begin{figure}[H]
    \centering
    \includegraphics[width=0.95\textwidth]{Figs/irregular_CHO_GS46_F_all_pl40.png}
    \caption{\textcolor{blue}{EnKF state estimation of Cell Line B Feed U plus 40\% under an irregular measurement schedule of alternating 48-hour and 72-hour sampling intervals. EnKF trajectories (solid lines) are compared against the irregularly spaced experimental measurements (triangles, $\pm$standard deviation). All eight state variables are tracked accurately despite the reduced and non-uniform measurement frequency.}}
    \label{fig:irregular}
\end{figure}

